@ID: OW24D1P1I
@BIDANG: Analisis Real

Jika $\mathcal I$ adalah koleksi interval buka pada $\mathbb R$ yang didefinisikan oleh

$$
\mathcal I=\left\{I_n:I_n=\left(1,1+\frac1n\right),\ n\in\mathbb N\right\},
$$

tentukan

$$
\bigcup_{I_n\in\mathcal I}I_n
\qquad\text{dan}\qquad
\bigcap_{I_n\in\mathcal I}I_n.
$$

@ID: OW24D1P2I
@BIDANG: Analisis Real

Untuk setiap $n\in\mathbb N$, definisikan fungsi $f_n:[-1,1]\to\mathbb R$ dengan

$$
f_n(x)=\cos(n\arccos x),\qquad x\in[-1,1].
$$

Jika $m,n\in\mathbb N$ dan $m\ne n$, tentukan nilai

$$
\int_{-1}^{1}\frac{f_n(x)f_m(x)}{\sqrt{1-x^2}}\,dx.
$$

@ID: OW24D1P3I
@BIDANG: Struktur Aljabar

Diberikan $G$ suatu grup siklis dengan orde $2024$. Banyak elemen di $G$ yang berorde ganjil adalah $\ldots$

@ID: OW24D1P4I
@BIDANG: Struktur Aljabar

Jika $a,b,c,d$ merupakan bilangan bulat sehingga

$$
f(x)=x^4+ax^3+bx^2+cx+d
$$

mempunyai akar $\sqrt3-\sqrt5$, maka tentukan nilai

$$
a+b+c+d.
$$

@ID: OW24D1P5I
@BIDANG: Kombinatorika

Dalam sebuah kotak terdapat kelereng berwarna biru, hijau, merah, kuning, dan abu-abu dengan masing-masing warna terdapat $9$ kelereng. Banyak cara mengambil $9$ kelereng dari kotak sehingga setiap warna paling sedikit terambil satu kali adalah $\ldots$

@ID: OW24D1P1U
@BIDANG: Analisis Real

Misalkan $a$ dan $b$ bilangan real positif yang memenuhi

$$
\sqrt b<a<2\sqrt b.
$$

Barisan bilangan real $(x_n)$ didefinisikan oleh $x_0\ge0$ dan

$$
x_n=\frac{ax_{n-1}+b}{x_{n-1}+a},
\qquad \forall n\in\mathbb N.
$$

Apakah $\displaystyle\lim_{n\to\infty}x_n$ ada? Jika ya, tentukan nilai limitnya.

@ID: OW24D1P2U
@BIDANG: Analisis Real

Diberikan fungsi kontinu $f:[a,b]\to\mathbb R$ dengan sifat bahwa untuk setiap $x\in[a,b]$, terdapat $p\in[a,b]$ sehingga

$$
|f(p)|\le\frac{2023}{2024}|f(x)|.
$$

Buktikan bahwa terdapat $c\in[a,b]$ sehingga $f(c)=0$.

@ID: OW24D1P3U
@BIDANG: Struktur Aljabar

Diketahui

$$
R=\mathbb Z_{1013}\times\mathbb Z_{1013}
$$

merupakan ring terhadap operasi penjumlahan dan perkalian berikut:

$$
(a,b)+(c,d)
=
\bigl((a+c)\bmod1013,\,(b+d)\bmod1013\bigr),
$$

$$
(a,b)\cdot(c,d)
=
\bigl(ac\bmod1013,\,(ad+bc+bd)\bmod1013\bigr),
$$

untuk setiap $(a,b),(c,d)\in R$.

Buktikan bahwa terdapat tepat sebanyak $2024$ elemen taknol di $R$ yang merupakan pembagi nol.

(Catatan: $1013$ merupakan bilangan prima.)

@ID: OW24D1P4U
@BIDANG: Struktur Aljabar

Suatu grup hingga $(G,\star)$ berorde $n$ dikatakan rapi jika terdapat $n$ unsur berbeda
$g_1,g_2,\ldots,g_n$ dari $G$ sehingga

$$
G=\{g_1\star g_2,\;g_2\star g_3,\;\ldots,\;g_{n-1}\star g_n,\;g_n\star g_1\}.
$$

a) Tunjukkan bahwa $(\mathbb Z_7,+)$ rapi.

b) Buktikan bahwa untuk setiap $n$ genap, $(\mathbb Z_n,+)$ tidak rapi.

@ID: OW24D1P5U
@BIDANG: Kombinatorika

Tunjukkan bahwa jika $n+1$ bilangan bulat berbeda diambil dari himpunan

$$
\{1,2,\ldots,kn\},
$$

maka selalu ada dua bilangan yang selisihnya paling banyak $k-1$.


@ID: OW24D2P1I
@BIDANG: Analisis Kompleks

Banyak bilangan kompleks $z$ yang memenuhi

$$
z^{20}=1
\quad\text{dan}\quad
z^{24}\in\mathbb R
$$

adalah $\ldots$

@ID: OW24D2P2I
@BIDANG: Analisis Kompleks

Jika diberikan fungsi kompleks

$$
f(z)=\overline{z}e^{-|z|^2},
$$

maka nilai $f'(i)$ adalah $\ldots$

@ID: OW24D2P3I
@BIDANG: Aljabar Linear

Misal $P_4$ adalah ruang polinomial berkoefisien real berderajat paling tinggi $4$.

Didefinisikan pemetaan $T_k:P_4\to P_4$ dengan

$$
T_k(a_0+a_1x+a_2x^2+a_3x^3+a_4x^4)
=
a_0+a_1^kx^2+a_2x^4.
$$

Banyak pasangan bilangan bulat $(k,m)$ dengan $k,m\in\{1,2,3,4\}$ sehingga $T_k^m$ merupakan pemetaan linear adalah $\ldots$

(Catatan: $T_k^m=\underbrace{T_k\circ T_k\circ\cdots\circ T_k}_{m\text{ kali}}$.)

@ID: OW24D2P4I
@BIDANG: Aljabar Linear

Misalkan $e_1,e_2,\ldots,e_{10}$ basis baku dari ruang vektor $\mathbb R^{10}$. Tinjau dua himpunan

$$
U=\{e_1,\;e_2+e_3,\;e_4+e_5+e_6,\;e_7+e_8+e_9+e_{10}\}
$$

dan

$$
V=\{e_1+e_2,\;e_3+e_4,\;e_5+e_6,\;e_7+e_8,\;e_9+e_{10}\}.
$$

Dimensi terkecil dari subruang yang memuat $U$ dan $V$ adalah $\ldots$

@ID: OW24D2P5I
@BIDANG: Kombinatorika

Dalam kejuaraan catur yang diikuti oleh $10$ peserta, setiap peserta bertanding dengan peserta lain tepat satu kali. Peserta yang menang, kalah, dan seri di setiap pertandingan berturut-turut diberikan skor $2$, $0$, dan $1$. Jika total skor setiap peserta di akhir kejuaraan berbeda-beda, maka maksimal banyak pertandingan yang mungkin dimenangkan oleh peserta dengan skor terendah adalah $\ldots$


@ID: OW24D2P1U
@BIDANG: Analisis Kompleks

Untuk setiap bilangan kompleks $z$ dengan $|z|=1$, tunjukkan bahwa

$$
1\le |1+z|+|1+2z|\le5.
$$

@ID: OW24D2P2U
@BIDANG: Analisis Kompleks

Tentukan daerah $D$ di bidang kompleks sehingga untuk setiap $z\in D$,

$$
\lim_{\substack{|z|\to\infty\\ z\in D}}e^z
$$

ada.

@ID: OW24D2P3U
@BIDANG: Aljabar Linear

Diberikan matriks

$$
A=
\begin{pmatrix}
0&a&0\\
a&b&a\\
0&a&0
\end{pmatrix},
$$

dengan $a,b\in\mathbb R$ dan $a\ne0$. Buktikan bahwa matriks $A$ dapat didiagonalkan.


@ID: OW24D2P4U
@BIDANG: Aljabar Linear

Diberikan matriks

$$
A=
\begin{pmatrix}
20&24\\
p&q
\end{pmatrix},
$$

dengan $p,q\in\mathbb R$.

Selidiki apakah terdapat pasangan bilangan real $(p,q)$ sehingga untuk suatu $b\in\mathbb R^2$, persamaan

$$
A^2x=Ab
$$

tidak memiliki solusi $x\in\mathbb R^2$. Jika ada, tentukan semua pasangan $(p,q)$ yang memenuhi.

@ID: OW24D2P5U
@BIDANG: Aljabar Linear

Jika koefisien $a$ dan $b$ dari persamaan garis lurus

$$
ax+by=0
$$

adalah dua bilangan berbeda dari himpunan $\{0,1,2,3,6,7\}$, tentukan banyaknya garis lurus berbeda yang dapat dibentuk.